\section{Discussion}
\label{sec:discussion}

\subsection{What This Study Shows, and What It Does Not}

CICADAS is an end-to-end modeling, estimation, and simulation framework for emulating randomized
trials of dynamic treatment policies from longitudinal observational data. In this paper we
specified the framework, implemented it in MATLAB and Python, and evaluated it entirely by internal,
simulation-based validation.

Within that evaluation, three things were shown. First, in synthetic ICU cohorts with confounding by
indication severe enough to eliminate---and in a substantial fraction of replications reverse---the
apparent treatment benefit, the parametric g-formula recovers the ground-truth policy contrast
closely, while naive Kaplan--Meier, stabilized IPTW, and a baseline-weighted marginal structural
model do not. Second, the sequential estimation procedure---natural history from untreated patients,
PK/PD parameters from dose-switching patients, mortality from the full cohort---recovers the
generating parameters accurately, and the corrected two-stage estimator approaches oracle
performance. Third, the estimator's behavior degrades in characterizable ways under injected
unmeasured confounding, EEG measurement error, and misspecified censoring.

None of this constitutes evidence about anti-seizure treatment in patients. Every number reported
here is a property of a data-generating process that we specified, recovered by an estimator that we
also specified; the two share structural assumptions, as Section~\ref{sec:misspec} discusses. The
appropriate reading is that the framework is \emph{methodologically feasible}---the estimation
problem is solvable, the computations are tractable at realistic cohort sizes, and the failure modes
are identifiable---and that its clinical utility is entirely undetermined.

\subsection{Methodological Contributions}

CICADAS addresses three challenges in ICU causal inference through integrated solutions: mechanistic
modeling captures underlying disease biology and drug effects, sequential parameter estimation
leverages natural treatment variation in clinical practice, and the parametric g-formula simulates
counterfactual outcomes under alternative policies. The corrected two-stage estimation approach
achieves near-oracle accuracy for key PK/PD parameters, while the g-formula removes the substantial
bias present in standard survival analyses of the same simulated data.

Causal effect recovery depends on accurate PK/PD parameter estimation, since these parameters
determine the mapping from treatment decisions to disease suppression and downstream mortality. In
these simulations, sufficient dosing variation was necessary for stable identification, whereas data
from constant closed-loop control provided insufficient dynamic range for reliable
estimation---motivating our use of dose-switching patients for parameter estimation. This is a
design insight that transfers directly to any real-data application of this class of method.

\subsection{Model Misspecification and the Shared-Structure Problem}
\label{sec:misspec}

The most important limitation of a simulation study of this design is that the data-generating
process and the estimation models are not independent. Both use a one-compartment pharmacokinetic
approximation; both represent pharmacodynamic effect as multiplicative suppression of disease burden
through a Hill function; both express mortality as a discrete-time hazard depending on cumulative
disease and cumulative drug burden through the same functional form; and both draw natural history
from the same parametric family. When the estimator is handed data generated under exactly the
structure it assumes, recovering the truth demonstrates that the estimator is consistent and that
the numerical implementation is correct. It does not demonstrate that the estimator is robust,
because nothing was misspecified.

We therefore regard the recovery results as characterizing an \emph{upper bound} on achievable
performance rather than an expected performance in practice. The sensitivity analyses in
Section~\ref{sec:robustness} relax this in three specific directions---an unmeasured confounder
entering both treatment assignment and mortality, additive measurement error on the EEG severity
proxy, and alternative censoring mechanisms---and the framework's behavior under each is
informative. But three named departures are not robustness to misspecification in general. In
particular, we did not test: the wrong number of PK compartments; non-monotone or saturating
dose--response departing from the Hill form; time-varying PK/PD parameters (which occur clinically
with changing renal and hepatic function, and with autoinduction); drug--drug interaction in
patients receiving multiple agents; or a mortality hazard depending on the \emph{pattern} of
exposure rather than its cumulative total. Any of these is plausible in real ICU data, and each
would bias the estimator in a direction that our current experiments cannot predict.

A distinct concern is transportability. Even a correctly specified model, fit in one cohort,
estimates a quantity specific to that cohort's case mix, drug formulary, monitoring intensity, and
---critically in critical care---its local practice around withdrawal of life-sustaining therapy.
Between-center variation in these is substantial, and it enters the estimand rather than merely the
variance. Estimates from a single center should not be assumed to transport, and a multicenter
analysis should report between-center heterogeneity rather than pooling over it.

\subsection{Conditions Required to Apply CICADAS to Real Data}
\label{sec:conditions}

The framework is not applicable to an arbitrary longitudinal ICU dataset. Six conditions must hold,
and each is checkable before any causal estimate is produced. We state them explicitly because a
framework whose preconditions are left implicit invites misapplication.

\begin{enumerate}\setlength{\itemsep}{3pt}
\item \textbf{Treatment overlap over the policy range of interest.} Every policy compared must have
support: for each patient stratum and each time point, the observed data must contain patients
receiving treatment intensities consistent with the policy being evaluated. In practice this means
examining the empirical distribution of dose conditional on severity and confirming it is
non-degenerate over the range of $\theta$ being swept. Where it is not, the corresponding portion of
the policy curve is extrapolation from the parametric model and should be reported as such rather
than plotted alongside supported regions (see Section~\ref{sec:structural-limits}).

\item \textbf{Measured confounding sufficient for sequential exchangeability.} The covariates
recorded must capture what actually drives dose escalation---not only severity and vital signs, but
sedation goals, neurologic examination findings, and goals-of-care status. The last of these is
routinely absent from structured data and is plausibly a strong confounder, since a decision to
limit escalation and a high mortality risk share a common cause. Negative-control
outcomes~\cite{lipsitch2010negative} and multiple propensity specifications should accompany any
real-data estimate.

\item \textbf{Dosing variation adequate to identify PK/PD parameters.} Section~\ref{sec:pkpd-results}
shows that patients under stable closed-loop control carry little information about drug effect:
identification requires within-patient variation in exposure. A cohort in which dosing is uniformly
protocolized is, for this purpose, uninformative---the framework requires the very practice
variation that motivates it. We recommend reporting the proportion of person-time spent with
$x_t/C_{50}$ in a range that supports identification, as a design diagnostic prior to estimation.

\item \textbf{A severity readout of adequate fidelity.} Our measurement-error experiments indicate
that the two-stage PK/PD estimator fails to converge on a substantial fraction of cohorts once noise
on the severity signal exceeds $\sigma \approx 0.05$. This is a concrete, quantitative precondition
on the EEG quantification pipeline, and it should be verified against that pipeline's known error
characteristics before analysis. Section~\ref{sec:eeg-readout} addresses the harder question of
whether a scalar readout is the right object at all.

\item \textbf{A characterized censoring mechanism.} Our censoring experiments show bias growing to
approximately $+17$ percentage points under strongly outcome-dependent censoring. In ICU data, loss
to follow-up is generated by transfer, monitoring discontinuation, and withdrawal of care---the last
of which is not plausibly independent of prognosis. The mechanism must be characterized and modeled,
not assumed.

\item \textbf{Subgroups represented in sufficient number.} Policy comparisons within a stratum
require that the stratum contain enough patients, with enough treatment variation, to fit the
component models. Strata failing this should be reported as indeterminate rather than as showing no
effect.
\end{enumerate}

\subsection{What the Framework Structurally Cannot Do}
\label{sec:structural-limits}

Two limits are structural rather than statistical, and no amount of data or computation removes
them.

The first bounds the credible policy space. CICADAS produces empirically defensible estimates only
for policies with support in the observed data; policies outside that range are extrapolations from
the fitted parametric models, and their apparent precision is an artifact of the parametric form
rather than evidence. This has a direct consequence for the threshold optimization presented in
Section~\ref{sec:policy-space}. Sweeping $\theta$ toward zero corresponds to a degree of suppression
that may simply not occur in routine practice, and in real data the corresponding portion of the
curve would be model-based conjecture. The honest output of a real-data application is therefore two
maps rather than one: a map of estimated treatment-effect heterogeneity, and a map of where the data
can speak at all. Regions where the observational data are structurally silent should be flagged for
prospective study, never presented as null findings.

The second limit is more fundamental and deserves to be stated plainly. Because the framework learns
the effects of treatment from variation in treatment that has already occurred, it can only ever
interrogate the space of therapies, mechanisms, and timing strategies embodied in current practice.
It can find the best policy within that space; it cannot find a policy outside it. A genuinely novel
agent, a mechanism no current drug targets, or a timing strategy nobody has attempted---the
disruptive interventions that produce real advances in clinical care---are invisible to it by
construction. This is not a deficiency to be engineered away; it is the definitional boundary
between learning from observation and learning from experiment. Randomized trials remain the only
instrument capable of evaluating interventions that have never been tried, and the appropriate
relationship between CICADAS and trials is therefore complementary: the framework is a tool for
allocating trial resources efficiently within known therapeutic space, not a substitute for the
trials that extend it.

\subsection{The EEG-Derived Treatment Target}
\label{sec:eeg-readout}

The framework assumes the availability of a scalar severity readout $L_t$ that can be tracked
continuously and titrated against a threshold. This assumption is doing more work than it appears
to, and we regard it as the most substantial obstacle to real-world deployment.

Clinically distinct electrographic states are not interchangeable instances of a single quantity.
Several brief focal seizures and one prolonged ictal episode may integrate to the same burden while
carrying different prognostic implications; focal and generalized activity differ in mechanism and
in likely response to systemic suppression; and identical ictal burden on a continuous, reactive
background versus a burst-suppressed or attenuated background reflects different underlying brain
states and, plausibly, different treatment benefit. There is also an irreducible measurement problem
beneath the modeling one. Expert electroencephalographers identifying seizures and rhythmic and
periodic patterns---precisely the categories a severity readout would have to be built
from---agree only moderately with one another~\cite{jing2023interrater}. A scalar target inherits
that irreducible label noise before any modeling assumption is made. Collapsing all of this to one
number then imposes a further substantive and currently unvalidated assumption: that harm is a
monotone function of a scalar summary, and that the summary is sufficient for the treatment
decision.

There is a practical problem alongside the conceptual one. Retrospective ICU datasets frequently
lack granular quantitative EEG; the recorded product is often a narrative report at daily resolution
rather than a time series at the resolution titration requires. Where raw EEG is available, deriving
a continuous target requires an automated quantification pipeline whose error characteristics are
known and whose output has been validated against outcomes---and, per
Section~\ref{sec:conditions}, whose noise falls below the threshold at which our estimation
procedure destabilizes.

The consequence is that optimizing $\theta$ is mathematically well posed inside the simulation while
being potentially uninterpretable outside it. A real deployment must therefore either commit to a
specific, prospectively validated scalar readout and defend that choice on its own evidence, or
generalize the framework so that severity is a vector or a pattern type rather than a scalar. The
latter is a substantial extension: it changes the controller, the identification argument, and the
policy space simultaneously. We regard resolving this as a prerequisite for, rather than a
refinement of, clinical application.

\subsection{Clinical Heterogeneity and the Mismatch with the Simulation}

The present simulation treats all critically ill patients with electrographic seizures as a single
population. In reality, causes of ICU seizures and the underlying pathophysiology are heterogeneous,
and the likelihood that anti-seizure treatment changes outcomes plausibly differs across etiologies.
Clinically important categories include: (i) post-anoxic encephalopathy, where widespread ischemic
injury from cardiac arrest may render treatment ineffective regardless of seizure
burden~\cite{rossetti2005refractory,ruijter2022treating}; (ii) primary structural injury
(large-artery ischemic stroke, intracerebral hemorrhage, traumatic brain injury, brain tumor), where
seizures arise from focal cortical damage and the benefit of aggressive suppression may depend on
lesion type and location~\cite{claassen2004detection,vespa2005continuous}; (iii) systemic,
metabolic, and septic encephalopathy, where seizures are often transient and reverse with treatment
of the underlying condition, so that aggressive anti-seizure therapy may add toxicity without
benefit~\cite{sutter2016status}; and (iv) primary epileptic populations with non-convulsive status
epilepticus, where conventional anti-seizure protocols are more likely to be
effective~\cite{trinka2015definition,husain2018randomized}. The causal structure---which covariates
confound, which mediate, and which modify treatment effects---likely differs across these groups.

We note explicitly that these etiologic distinctions are \textbf{not represented in the present
simulation}. The heterogeneity in our data-generating process is parametric---patients differ in
susceptibility to cumulative disease harm and to cumulative drug harm---rather than mechanistic.
Parametric heterogeneity of this kind is sufficient to exercise the machinery of stratified policy
estimation, but it is not the same thing as the clinical heterogeneity we invoke to motivate the
framework's relevance, and it should not be read as evidence that the framework handles that
heterogeneity. A simulation that did so would need distinct disease-dynamics and benefit--harm
structures per etiology, including etiologies in which treatment is futile at any dose, and would
need to test whether the estimator recovers effects separately within strata of markedly unequal
size. We regard building that simulation as a necessary step before real-data deployment, and it is
not attempted here.

Similarly, on the harm side of the trade-off, the present simulation represents treatment harm only
as a concentration-dependent contribution to the mortality hazard. A complete representation in a
real clinical deployment would separate concrete harms that clinicians weigh in practice:
sedation-related hypotension and the need for vasopressor support, prolonged mechanical ventilation
and the associated risk of ventilator-associated pneumonia, delayed weaning and ICU length of stay,
delayed neuroprognostication in post-arrest patients, and drug-specific toxicities (e.g., propofol
infusion syndrome~\cite{kam2007propofol}, thrombocytopenia with valproate, hepatic failure with other
agents~\cite{perucca2012adverse}). The framework's outcome model can be extended to treat each of
these as a distinct competing or intermediate outcome, and we regard this as a priority in the
real-data application.

\subsection{Limitations}

Beyond the structural issues above, our simulations assume simplified patient heterogeneity and
one-compartment pharmacokinetics, whereas real applications may require additional covariates and
more complex models. The identification argument assumes no unmeasured confounding of the
treatment--outcome relationship at each time point; Section~\ref{sec:robustness} quantifies the
bias induced by violations of a specific parametric form, but cannot establish that the assumption
holds in any real dataset. Finally, our benchmarking comparisons are subject to the caveat stated in
Section~\ref{sec:benchmark}: our g-formula implementation is correctly specified with respect to the
generating process while its competitors are off-the-shelf specifications, so the ordering we
observe should not be read as a general ranking of methods.

\subsection{External Validation}
\label{sec:external-validation}

The necessary next step is external validation in real ICU EEG data, and we describe it here
strictly as future work that does not support any claim in the present manuscript. We are assembling
a multicenter cohort of ICU patients undergoing continuous EEG monitoring for this purpose. That
analysis will be the first test of whether the preconditions in Section~\ref{sec:conditions} are
satisfiable in practice---whether overlap is adequate, whether the measured covariates suffice,
whether dosing variation identifies PK/PD parameters, and whether a severity readout of the required
fidelity can be constructed. It is entirely possible that they are not, and a negative answer on any
of them would be an informative result about the limits of this class of method.

Only if those preconditions are met does the question of clinical utility arise, and even then the
framework's output would be a set of hypotheses about which patients and which treatment intensities
merit prospective evaluation---not treatment recommendations. Establishing clinical utility would
require a subsequent trial testing a CICADAS-derived strategy against usual care. Until such a trial
is conducted, the framework's value is methodological.
